% ============================================================
% Round 08: Dr.GRPO
% ============================================================

\subsection{Round 08：Dr.GRPO}

\subsubsection{本轮问题}

\begin{conceptbox}
\textbf{学习目标：} 理解 Dr.GRPO 为什么认为原始 GRPO / DAPO 仍然有优化偏置，以及它如何通过去掉 response length normalization 和 reward std scaling 来提高 token efficiency。\\
\textbf{主线位置：}
\[
\text{GRPO}
\to
\text{DAPO}
\to
\text{Dr.GRPO}
\to
\text{GSPO}
\]
\end{conceptbox}

\begin{conceptbox}
\textbf{Highlight：Dr.GRPO}\\
\textbf{论文：} \emph{Understanding R1-Zero-Like Training: A Critical Perspective}\\
\textbf{arXiv：} \href{https://arxiv.org/abs/2503.20783}{2503.20783}，submitted 2025-03-26\\
\textbf{OpenReview：} \href{https://openreview.net/forum?id=5PAF7PAY2Y}{COLM 2025 submission / revision page}\\
\textbf{Code：} \href{https://github.com/sail-sg/understand-r1-zero}{sail-sg/understand-r1-zero}\\
\textbf{关键词：} R1-Zero-like training，GRPO bias，length bias，difficulty bias，critic-free RL，token efficiency\\
\textbf{核心定位：} 批判性分析 R1-Zero 类训练中的优化偏置，提出 Dr.GRPO 来减少错误回答变长和 token 浪费。
\end{conceptbox}

\subsubsection{Dr.GRPO 想修什么}

DAPO 已经在 GRPO 上做了很多工程改进：

\[
\text{Dynamic Sampling}
\quad
\text{Clip-Higher}
\quad
\text{Token-Level Loss}
\quad
\text{Overlong Shaping}
\]

但 Dr.GRPO 关注的问题更“数学化”：它认为原始 GRPO objective 里有两个容易被忽略的优化偏置：

\begin{enumerate}[leftmargin=1.5em]
  \item \textbf{response length bias：} 由每条回答自己的长度归一化 $1/|o_i|$ 引入；
  \item \textbf{difficulty / std scaling bias：} 由组内 reward 标准差归一化 $1/\operatorname{std}(\{r_j\})$ 引入。
\end{enumerate}

\begin{summarybox}
\textbf{一句话：} Dr.GRPO = Group Relative Policy Optimization Done Right。它的核心不是加新技巧，而是删掉会扭曲梯度权重的归一化项。
\end{summarybox}

\subsubsection{先看原始 GRPO 的两个归一化}

原始 GRPO 常见 loss 写法是：

\[
\mathcal{L}_\text{GRPO}
=
-
\frac{1}{G}
\sum_{i=1}^{G}
\frac{1}{|o_i|}
\sum_{t=1}^{|o_i|}
l_{i,t}
\]

其中：

\[
l_{i,t}
\approx
\rho_{i,t}(\theta)\hat{A}_i
\]

而 group advantage 通常是：

\[
\hat{A}_i
=
\frac{
r_i-\operatorname{mean}(\{r_j\}_{j=1}^{G})
}{
\operatorname{std}(\{r_j\}_{j=1}^{G})+\epsilon
}
\]

所以这里有两个归一化：

\begin{center}
{\small
\begin{tabular}{p{4cm}p{8.8cm}}
\hline
\textbf{归一化项} & \textbf{作用} \\
\hline
$1/|o_i|$
&
每条回答内部按自己的长度平均 token loss。
\\
$1/\operatorname{std}(\{r_j\})$
&
把同一个 prompt 的 group reward 标准化到近似单位方差。
\\
\hline
\end{tabular}
}
\end{center}

Dr.GRPO 的核心观点是：这两个项看起来像正常化技巧，但会改变不同样本的真实优化权重。

\subsubsection{问题一：Response Length Bias}

先看长度项：

\[
\frac{1}{|o_i|}
\sum_{t=1}^{|o_i|}
l_{i,t}
\]

这表示每条回答不管多长，都会先变成一个平均值。

如果两个回答的 advantage 相同：

\[
\hat{A}_1 = \hat{A}_2
\]

但长度不同：

\[
|o_1| = 100,
\qquad
|o_2| = 1000
\]

那么原始 GRPO 会让两条回答的总 loss 权重差不多，而不是让长回答里的 1000 个 token 都按同样强度参与更新。

\begin{warnbox}
\textbf{关键问题：} 这种按 response 自身长度平均的写法，会让回答长度参与决定梯度权重。Dr.GRPO 认为这会造成 response-level length bias。
\end{warnbox}

在 reasoning RL 中，这个偏置特别麻烦，因为模型的回答长度本身也在训练中变化。也就是说，模型不只是学习“怎么答对”，还可能被 objective 暗中鼓励某些长度行为。

论文特别指出，原始 GRPO 可能使回答长度持续增长，尤其是错误回答变长。这会造成：

\begin{itemize}[leftmargin=1.5em]
  \item token efficiency 下降；
  \item 错误回答里出现更多无效思考；
  \item 表面上像 long-CoT emergent behavior，但其中一部分其实来自优化偏置。
\end{itemize}

\subsubsection{Dr.GRPO 怎么修长度偏置}

Dr.GRPO 不再用每条回答自己的长度 $|o_i|$ 做归一化，而是用一个固定常数 $L$。通常 $L$ 可以取最大生成长度：

\[
L = T_\text{max}
\]

Dr.GRPO loss 写成：

\[
\mathcal{L}_\text{Dr.GRPO}
=
-
\frac{1}{LG}
\sum_{i=1}^{G}
\sum_{t=1}^{|o_i|}
l_{i,t}
\]

对比一下：

\[
\mathcal{L}_\text{GRPO}
=
-
\frac{1}{G}
\sum_i
\frac{1}{|o_i|}
\sum_t l_{i,t}
\]

\[
\mathcal{L}_\text{DAPO}
=
-
\frac{1}{\sum_i |o_i|}
\sum_i
\sum_t l_{i,t}
\]

\[
\mathcal{L}_\text{Dr.GRPO}
=
-
\frac{1}{LG}
\sum_i
\sum_t l_{i,t}
\]

\begin{intubox}
\textbf{直觉：} 原始 GRPO 是“每条回答用自己的长度归一化”；DAPO 是“用当前 batch 的总 token 数归一化”；Dr.GRPO 是“用固定最大长度归一化”。固定常数不随模型生成长度变化，因此更不容易引入长度偏置。
\end{intubox}

\subsubsection{问题二：Reward Std Scaling Bias}

标准 GRPO 的 advantage 常写成：

\[
\hat{A}_i
=
\frac{
r_i-\mu_\text{group}
}{
\sigma_\text{group}+\epsilon
}
\]

这里：

\[
\mu_\text{group}
=
\operatorname{mean}(\{r_j\})
\]

\[
\sigma_\text{group}
=
\operatorname{std}(\{r_j\})
\]

均值中心化是合理的，它让同组回答相对比较：

\[
r_i-\mu_\text{group}
\]

但再除以组内标准差，会改变不同 prompt 的权重。

举例说，同样是二值 reward，如果一个 group 的正确率接近一半：

\[
[0,1,0,1]
\]

它的标准差较大。

如果一个 group 只有一个样本和其他不一样：

\[
[0,0,0,1]
\]

它的标准差较小。除以更小的 $\sigma$ 会放大这类 group 的梯度。

\begin{warnbox}
\textbf{关键问题：} reward std scaling 不只是“数值稳定”，它会让不同难度、不同 reward 分布的 prompt 获得不同训练权重。Dr.GRPO 把这称为一种 difficulty / question-level bias。
\end{warnbox}

\subsubsection{Dr.GRPO 怎么修 Std Scaling Bias}

Dr.GRPO 建议去掉组内 std scaling，只保留均值 baseline：

\[
\hat{A}_i
=
r_i-\operatorname{mean}(\{r_j\}_{j=1}^{G})
\]

也就是说，从：

\[
\frac{r_i-\mu_\text{group}}{\sigma_\text{group}+\epsilon}
\]

改成：

\[
r_i-\mu_\text{group}
\]

\begin{intubox}
\textbf{直觉：} 均值 baseline 负责告诉模型“比同组平均好还是差”；std scaling 会额外改变这个 prompt 的整体权重。Dr.GRPO 认为后者会带来不想要的难度偏置。
\end{intubox}

\subsubsection{Dr.GRPO 的完整改法}

因此 Dr.GRPO 的两个主要修改是：

\begin{enumerate}[leftmargin=1.5em]
  \item 不用每条回答的 $|o_i|$ 做归一化，改用固定常数 $L$；
  \item 不用组内 reward std scaling，advantage 用 $r_i-\mu_\text{group}$。
\end{enumerate}

可以把它简化写成：

\[
\mathcal{L}_\text{Dr.GRPO}
=
-
\frac{1}{LG}
\sum_{i=1}^{G}
\sum_{t=1}^{|o_i|}
\rho_{i,t}(\theta)
\left(
r_i-\mu_\text{group}
\right)
\]

实际实现中仍然可以保留 PPO-style clip、old policy ratio 等机制；Dr.GRPO 的核心是在 loss normalization 和 advantage scaling 上去偏。

\subsubsection{GRPO / DAPO / Dr.GRPO 对比}

\begin{center}
{\small
\begin{tabular}{p{3cm}p{3.2cm}p{3.2cm}p{3.2cm}}
\hline
 & \textbf{GRPO} & \textbf{DAPO} & \textbf{Dr.GRPO} \\
\hline
长度归一化
&
$1/|o_i|$
&
$1/\sum_i |o_i|$
&
$1/(LG)$
\\
是否依赖生成长度
&
依赖单条长度
&
依赖 batch 总长度
&
不依赖，使用固定常数
\\
reward std scaling
&
通常使用
&
通常沿用
&
建议去掉
\\
主要目标
&
critic-free online RL
&
工程稳定性
&
去除优化偏置，提高 token efficiency
\\
\hline
\end{tabular}
}
\end{center}

\subsubsection{Dr.GRPO 的直觉}

Dr.GRPO 的思想可以很朴素地理解：

\[
\text{不要让回答长度决定训练权重}
\]

\[
\text{不要让 group reward 标准差决定 prompt 权重}
\]

也就是说，模型应该因为答案质量被强化或惩罚，而不是因为这个回答刚好长、这个 group 的 reward 方差刚好小而被额外放大或缩小。

\begin{summarybox}
\textbf{一句话：} GRPO / DAPO 更像工程可用的 recipe；Dr.GRPO 更像对这些 recipe 的目标函数体检，指出哪些归一化会悄悄改变优化目标。
\end{summarybox}

\subsubsection{本轮最小结论}

\begin{summarybox}
\textbf{Dr.GRPO 的核心：}\\
去掉 response length normalization：
\[
\frac{1}{|o_i|}
\quad \to \quad
\frac{1}{L}
\]
去掉 reward std scaling：
\[
\frac{r_i-\mu}{\sigma}
\quad \to \quad
r_i-\mu
\]
它要解决的是 GRPO 中的 length bias 和 difficulty / std scaling bias。\\
核心效果：减少错误回答变长，提高 token efficiency，同时尽量保持 reasoning performance。
\end{summarybox}

\subsubsection{本轮检查题}

\begin{enumerate}[leftmargin=1.5em]
  \item 原始 GRPO 里的 $1/|o_i|$ 为什么可能引入长度偏置？
  \item DAPO 的 token-level loss 和 Dr.GRPO 的常数归一化有什么区别？
  \item reward std scaling 为什么可能带来 question-level difficulty bias？
  \item Dr.GRPO 为什么说自己更接近无偏 policy gradient？
  \item Dr.GRPO 和 DAPO 分别更像“理论修正”还是“工程 recipe”？
\end{enumerate}
